% ----------------------------------------------------------------------------------- %
% VII Col\'oquio de Matem\'atica da Regi\~ao Centro-Oeste
% UFGD - Dourados, MS - 08 a 12 de abril de 2024
% 
% O modelo fornece dois comandos para auxiliar ne declara\c{c}\~ao de figuras.
% 
% Uma \'unica figura:
% \figura{fator_de_escala}{nome_do_arquivo}{legenda_da_figura}
%
% Duas figuras lado a lado ocupando toda a largura da linha:
% \figuras{fator_de_escala}{arquivo_da_figura1}{arquivo_da_figura2}{legenda_da_figura}
%
% O comando \colunas organiza o conte\'udo do p\^oster em duas colunas.
% ----------------------------------------------------------------------------------- %

\documentclass{vii-coloquio-co}

\usepackage{lipsum}  % utilizado para gerar o texto do modelo

% declara\c{c}\~ao dos autores do trabalho
\titulo{T\'{\i}tulo bem criativo}
\autor{Primeiro Autor*, Segundo Autor**}
\email{*primeiro@autor.edu.br, **segundo@autor.edu.br}
\universidade{Universidade Federal da Grande Dourados}

\begin{document}

\cabecalho

\colunas{
    \section*{Resumo}
    \lipsum[1-2]
    
    \section{Primeira se\c{c}\~ao}
    \lipsum[1-2]\cite{texbook}

    \figura{1.0}{figura.png}{Figura com tamanho original.}

    \lipsum[3-5]
    
    \section{Segunda se\c{c}\~ao}
    \lipsum[1-2]\cite{latex2e}

    \figura{2.0}{figura.png}{Figura com o dobro do tamanho.}
    \figura{0.5}{figura.png}{Figura com metade do tamanho.}
    \figuras{figura.png}{figura.png}{Figuras lado a lado.}

    \lipsum[3-4]

    \section{Conclus\~ao}
    \lipsum[1]

    \section*{Agradecimentos}
    \lipsum[1]

    \bibliographystyle{plain}
    \bibliography{referencias}
}

\rodape
\end{document}
